\section{Refutations and the finite verification}\label{sec:certificates}
A SAT solver is used to find a refutation; a separate checker establishes
that the saved refutation is valid for the saved input. The distinction is
useful because the search may be complicated, whereas the inference rule
checked at each proof step is elementary. The complete computation used
here has three outputs: a covering list of branches, validated Boolean
inputs for those branches, and a checked refutation for each terminal input.

\subsection{What one proof step means}
Given clauses $D$, \emph{unit propagation} repeatedly sets the sole
unassigned literal of a clause to true when its other literals are false.
An empty clause or conflicting forced assignments gives a contradiction.
For example, from $(p\lor q)$ and $\neg p$, the temporary assumption
$\neg q$ forces $p$ and a conflict. Thus $D$ entails the unit clause $q$.

\begin{definition}
A clause $C=\ell_1\lor\cdots\lor\ell_k$ has the
\emph{reverse-unit-propagation} (RUP) property with respect to $D$ if unit
propagation on
\[
D\land\neg\ell_1\land\cdots\land\neg\ell_k
\]
produces a contradiction. For an empty $C$, there are no temporary
assumptions: $D$ itself must yield a unit-propagation contradiction.
\end{definition}

\begin{lemma}[Soundness of RUP checking]\label{lem:rup}
If $C$ is RUP with respect to $D$, then $D\models C$. Starting from a
formula $F$, a sequence of checked RUP additions that reaches the empty
clause proves $F$ unsatisfiable, provided every clause used in a check is
an original clause or a previously verified consequence of $F$.
Deleting previously verified clauses or retaining them preserves this
soundness statement.
\end{lemma}
\begin{proof}
Any satisfying assignment of $D\land\neg C$ would have to obey every
forced unit. It cannot obey conflicting units or satisfy an empty clause.
Hence $D\land\neg C$ is unsatisfiable, so $D\models C$.
Inductively, every retained clause is a consequence of $F$. A newly
accepted clause follows from retained consequences of $F$ and is
therefore also a consequence of $F$. Deleting some of those consequences,
or continuing to retain them, does not introduce an unjustified premise.
A checked empty clause gives $F\models\bot$.
\end{proof}

The saved files use the DRUP text format, with addition and deletion
records; see~\cite{drat} for the wider clausal-proof setting. Here every
addition must pass RUP: no general RAT rule is used. The full replay
reported below uses \texttt{FreshRup}, a C\# implementation that strictly
parses deletion records but ignores their deletion effect. The logical
database retains the original clauses and previously accepted
consequences; the implementation may represent forced units by permanent
root assignments. Retaining learned consequences costs memory but makes
the premise invariant especially simple. It is a custom checker, not a
run of DRAT-trim.

\subsection{The finite assertion}
\begin{proposition}[Finite exclusion by computation]\label{prop:finite}
For every $17\leq n\leq25$ and $a\in\calR_n$, the adopted evidence
contains either a checked RUP refutation of $F_a$, or checked RUP
refutations of all 51 formulas $F_{a,t}$, $t\in\calT$. The associated
coverage and input-preservation checks of
Proposition~\ref{prop:input} pass for these actual inputs.
\end{proposition}
\begin{proof}[Computational verification]
The count representatives and diagonal patterns are reconstructed from
their definitions and compared with the saved catalogues. The actual
clauses are checked against Section~\ref{sec:encoding}; valid input
assignments are extended and evaluated on every clause of each actual
counter type. The terminal list is then derived from the parent catalogue:
a direct parent contributes one leaf, and a refined parent contributes
exactly one child for each of the 51 patterns.

For every leaf, the audited CNF is identified by its path and SHA-256,
paired with the adopted proof, and passed directly to \texttt{FreshRup}.
Every addition, including the final empty clause, is checked from that
input. The full run on October 27, 2026 completed successfully for exactly
the selected leaf set. Table~\ref{tab:counts} gives its per-weight totals;
Section~\ref{sec:full-audit} describes the checks binding these totals to
the actual files. The sources, complete certificates and raw outputs are
identified in Appendices~\ref{app:reproduction} and~\ref{app:audit}.
\end{proof}

\begin{table}[htbp]
\centering\small
\caption{Complete classification and the October 27 full replay.
$N_n=|\calA_n|$, $R_n=|\calR_n|$, $B_n$ counts refined parents, and
$L_n$ counts terminal proofs. Each addition total includes the final
empty clause of each proof.}\label{tab:counts}\label{tab:fresh-audit}
\begin{tabular}{@{}rrrrrrr@{}}
\toprule
$n$&$N_n$&$R_n$&$B_n$&$51B_n$&$L_n$&RUP additions\\
\midrule
17&20&3&0&0&3&14,413\\
18&90&12&0&0&12&70,751\\
19&250&29&0&0&29&284,448\\
20&526&60&0&0&60&742,441\\
21&940&102&0&0&102&2,776,244\\
22&1,515&163&11&561&713&5,704,885\\
23&2,275&240&28&1,428&1,640&11,590,970\\
24&3,245&341&36&1,836&2,141&16,112,383\\
25&4,451&463&69&3,519&3,913&29,107,399\\
\midrule
Total&13,312&1,413&144&7,344&8,613&66,403,934\\
\bottomrule
\end{tabular}

\end{table}

Refinement replaces a parent by its 51 children, so
\begin{equation}\label{eq:leaf-count}
L_n=(R_n-B_n)+51B_n=R_n+50B_n.
\end{equation}
Consequently, the 1413 parents consist of 1269 direct parents and 144
refined parents. There are $144\cdot51=7344$ child inputs and
$1269+7344=8613$ terminal proofs. The input audit checks all
$1413+7344=8757$ inputs, including the 144 nonterminal parents needed to
validate the refinements. Those parents are not counted again as leaves.

\subsection{How the complete run is tied to the mathematical claim}\label{sec:full-audit}
The October 27 input auditor was a new Python standard-library
implementation. It imports neither the original encoder nor an earlier
auditor. It independently reconstructs the 400 base variables and 1595
base clauses, the count representatives, and the 51 patterns. For every
parent and child it performs the actual-clause checks described in
Section~\ref{sec:encoding}. Across all weights there are 14 actual
counter-component types, with targets $r=0,\ldots,13$. The complete
normalized clause list is part of a type's identity. All
\[
\sum_{r=0}^{13}\binom{16}{r}=65,399
\]
valid component inputs pass the extension check and direct clause
evaluation. This is exhaustive over the actual types, not random
sampling and not a claim about every encoding the library can generate.

The connection from the catalogue to the run is checked in the following
order; a failure in any step prevents a full-success report.
\begin{enumerate}
\item Reconstruct every parent representative. Require exactly one
adopted direct leaf or all 51 specified children for that parent.
\item Require unique leaf identifiers, CNF paths and proof paths. Match
every terminal CNF's hash to the input auditor's record, and check the
paired proof's size and hash against the selection ledger.
\item Run the same frozen checker on every selected pair. Require each
addition to pass and each proof to end in a checked empty clause.
\item Parse every final output row strictly. Match its identity, batch
position, addition count and empty-clause flag to the selected input.
Require all batch exits to succeed and the returned identity set to equal
the expected set.
\end{enumerate}
These checks are stronger than counting 8613 output rows: a duplicate
could otherwise hide a missing leaf. Hashes establish which bytes were
checked; the input and RUP checks establish their mathematical role.

All 8613 proofs passed, with 66,403,934 additions including 8613 final
empty clauses. The checker also parsed and ignored 42,830,405 deletion
records. All 17,226 terminal CNF/proof files, totaling 6,585,642,510 bytes,
matched their recorded sizes and hashes. The run used eight batches,
started at 00:58:43 and ended at 01:16:33 on October 27, 2026
(UTC+08:00), with a wall time of 1070.69 seconds. The sum of individual
leaf times was 6153.77 seconds; parallel wall time and summed task time
are different quantities. No old success cache or newly rerun SAT search
was used, and the checker source and executable were frozen throughout.

The original search used Glucose3 through PySAT. Some parent searches
ended with \texttt{UNKNOWN} or a conflict-budget limit and were later
replaced by complete 51-child refutations. Such statuses, interrupted
runs and unchecked \texttt{UNSAT} outputs do not exclude a branch.
The October 27 full run uses only \texttt{FreshRup}; active-deletion
variants and earlier fallback counts belong to the supplementary or
historical records in Appendix~\ref{app:audit}.
